%% ============================================================
%%  IEEE TWO-COLUMN CONFERENCE PAPER
%%  Topic : Etsy AI-Generated Image Detection
%%  Author: Indumanaswini Bhukya, Dublin City University
%%  Notes : Compile with pdflatex (or upload to Overleaf).
%%          Replace \framebox placeholders with real figures.
%% ============================================================
\documentclass[conference]{IEEEtran}
\IEEEoverridecommandlockouts

\usepackage{cite}
\usepackage{amsmath,amssymb,amsfonts}
\usepackage{graphicx}
\usepackage{textcomp}
\usepackage{xcolor}
\usepackage{booktabs}
\usepackage{array}
\usepackage{multirow}

\def\BibTeX{{\rm B\kern-.05em{\sc i\kern-.025em b}\kern-.08em
    T\kern-.1667em\lower.7ex\hbox{E}\kern-.125emX}}

\begin{document}

%% ── Title ────────────────────────────────────────────────────────────────────
\title{AI-Generated Image Detection for Etsy Product Listings:\\
       A Dual-Backbone Embedding Fusion Approach}

\author{
  \IEEEauthorblockN{Indumanaswini Bhukya}
  \IEEEauthorblockA{\textit{School of Computing} \\
    \textit{Dublin City University}\\
    Dublin, Ireland \\
    indumanaswini.bhukya2@mail.dcu.ie}
  \and
  \IEEEauthorblockN{Rajesh Suddala}
  \IEEEauthorblockA{\textit{School of Computing} \\
    \textit{Dublin City University}\\
    Dublin, Ireland \\
    rajesh.suddala2@mail.dcu.ie}
}

\maketitle

%% ── Abstract ─────────────────────────────────────────────────────────────────
\begin{abstract}
The proliferation of AI image generators such as Midjourney, DALL-E, and Stable
Diffusion has made it possible for sellers on e-commerce platforms like Etsy to
misrepresent listings using synthetic imagery instead of authentic product
photographs.  We present a binary classification pipeline that detects
AI-generated product images (class~1) versus authentic photographs (class~0).
Our approach combines two complementary pretrained backbones—ConvNeXt V2 Tiny
and Swin Transformer Tiny—and fuses their 256-dimensional penultimate-layer
embeddings into a 512-dimensional representation on which a CatBoost meta-learner
is trained.  A five-view test-time augmentation (TTA) strategy and a validation-set
threshold sweep further refine the decision boundary.  On a dataset of 4\,800
labelled images (80/20 train--validation split), the full pipeline achieves a
validation F1 score of \textbf{0.8771} at an optimal decision threshold of 0.68,
outperforming each backbone used in isolation.
\end{abstract}

\begin{IEEEkeywords}
AI-generated image detection, binary image classification, transfer learning,
ConvNeXt V2, Swin Transformer, embedding fusion, CatBoost, test-time augmentation,
e-commerce moderation
\end{IEEEkeywords}

%% ── 1. Introduction ──────────────────────────────────────────────────────────
\section{Introduction}

Etsy is a globally recognised two-sided online marketplace that connects millions
of independent sellers with buyers seeking handmade, vintage, and unique goods.
Recent advances in text-to-image generation—Midjourney, DALL-E~\cite{ramesh2022dalle},
and Stable Diffusion~\cite{rombach2022latent}—have lowered the barrier to creating
photorealistic synthetic images, enabling a small but growing number of sellers to
misrepresent their listings with AI-generated imagery.  At scale, manual review of
every listing image is infeasible; automated detection is therefore essential for
marketplace integrity.

The task is framed as supervised binary classification: given a product listing
image, predict whether it is AI-generated (label~1) or an authentic photograph
(label~0).  Performance is measured by the F1 score on a hidden test set, which
balances precision and recall, making it particularly appropriate for a
near-balanced dataset.

Our contributions are as follows.
\begin{itemize}
  \item We train two complementary pretrained backbones—ConvNeXt V2
        Tiny~\cite{woo2023convnextv2} and Swin Transformer
        Tiny~\cite{liu2021swin}—with a staged fine-tuning strategy and show
        that both achieve validation F1 ${\geq}$~0.84 independently.
  \item We fuse their 256-dimensional penultimate embeddings into a compact
        512-dimensional representation and benchmark three meta-learners
        (XGBoost~\cite{chen2016xgboost}, CatBoost~\cite{prokhorenkova2018catboost},
        and TabPFN) on this joint feature space.
  \item We demonstrate that five-view test-time augmentation combined with
        validation-set threshold optimisation improves the F1 score from 0.8639
        (CatBoost, fixed threshold~=~0.50) to \textbf{0.8771} (threshold~=~0.68).
  \item Gradient-weighted Class Activation Mapping (Grad-CAM)~\cite{selvaraju2017gradcam}
        is used to diagnose systematic failure modes, revealing that false
        positives arise primarily from misleading background textures, while
        false negatives are caused by severe lighting artefacts.
\end{itemize}

The remainder of the paper is structured as follows.  Section~\ref{sec:related}
reviews related work.  Section~\ref{sec:method} describes the full pipeline.
Section~\ref{sec:experiments} reports experimental results.
Section~\ref{sec:discussion} discusses findings and limitations.
Section~\ref{sec:conclusion} concludes.

%% ── 2. Related Work ──────────────────────────────────────────────────────────
\section{Related Work}
\label{sec:related}

\subsection{Synthetic Image Detection}

Early forensic detectors exploited compression artefacts and camera noise
patterns~\cite{farid2009image}.  Wang~\textit{et al.}~\cite{wang2020cnn} showed
that a ResNet-50 classifier trained on ProGAN images transfers surprisingly well
to other GAN families, indicating that CNNs learn generalised synthetic-texture
signatures.  Frank~\textit{et al.}~\cite{frank2020leveraging} demonstrated that
the frequency spectrum of GAN images contains characteristic peaks absent in
real photographs.  The emergence of diffusion-based generators~\cite{rombach2022latent}
has challenged older detectors; Corvi~\textit{et al.}~\cite{corvi2023detection}
report significant performance degradation when GAN-trained detectors face
diffusion-model outputs, motivating training on data from modern generators.
Ojha~\textit{et al.}~\cite{ojha2023towards} propose using CLIP~\cite{radford2021learning}
feature space for universal detection with strong cross-generator generalisation.

\subsection{Hierarchical and Convolutional Backbones}

ConvNeXt V2~\cite{woo2023convnextv2} extends the convolutional paradigm with
Fully Convolutional Masked Autoencoders (FCMAE), a self-supervised pretraining
objective that encourages the model to learn high-fidelity local texture
representations.  Swin Transformer~\cite{liu2021swin} introduces a hierarchical
shifted-window self-attention mechanism that captures both fine-grained local
detail and long-range global context at multiple resolutions, with memory
efficiency favourable to constrained hardware.  Both architectures have achieved
state-of-the-art results on ImageNet classification and downstream tasks.

\subsection{Ensemble Methods and Threshold Optimisation}

Gradient-boosted decision trees trained on deep features extracted from
pretrained networks constitute a well-established two-stage paradigm for
tabular and image tasks~\cite{chen2016xgboost,prokhorenkova2018catboost}.
Probability calibration and decision-threshold optimisation on a held-out
validation set are standard techniques for maximising F1 when the class
distribution is not severely skewed~\cite{provost2000wellbehaved}.

%% ── 3. Methodology ───────────────────────────────────────────────────────────
\section{Methodology}
\label{sec:method}

Figure~\ref{fig:pipeline} illustrates the end-to-end pipeline.  Data are loaded
from CSV manifests, processed through two separately trained backbones, and their
embeddings are fused before meta-learner training, TTA, and threshold optimisation.

\begin{figure}[t]
  \centering
  %% Replace with: \includegraphics[width=\columnwidth]{figs/pipeline.pdf}
  \framebox{\parbox{0.92\columnwidth}{\centering\vspace{1.2cm}%
    \texttt{train.csv / test.csv}
    $\;\rightarrow\;$ Augmentation
    $\;\rightarrow\;$ ConvNeXt V2 / Swin
    $\;\rightarrow\;$ 512-d Fusion
    $\;\rightarrow\;$ CatBoost
    $\;\rightarrow\;$ TTA + Threshold
    $\;\rightarrow\;$ \texttt{submission.csv}
    \vspace{1.2cm}}}
  \caption{End-to-end prediction pipeline.}
  \label{fig:pipeline}
\end{figure}

\subsection{Dataset and Exploratory Data Analysis}

The dataset comprises 4\,800 labelled training images and 2\,058 unlabelled test
images sourced from Etsy product listings.  Labels indicate whether an image is
AI-generated (1) or authentic (0).  The class distribution in the full training
set is 2\,485 authentic and 2\,315 AI-generated (imbalance ratio~=~1.07),
making it near-balanced.  A stratified 80/20 split yields 3\,840 training and
960 validation samples, preserving the class ratio in both subsets.

Prior to model training, an exploratory data analysis (EDA) step examines five
dimensions of the dataset to validate augmentation choices and surface potential
failure modes:
\begin{enumerate}
  \item \textbf{Class distribution} — a bar chart and pie chart confirm the
        near-balanced split (48.2\% AI-generated, 51.8\% authentic).
  \item \textbf{Sample image grids} — eight images per class are displayed to
        reveal that AI-generated images often exhibit smoother gradients and
        over-saturated colours compared with authentic photographs.
  \item \textbf{Image geometry} — a scatter plot of width vs.\ height and
        marginal histograms over 300 sampled images show that both classes share
        a broad resolution range; most images lie between 300 and 1200\,px on
        each side, with no systematic geometry gap between classes.
  \item \textbf{Per-class colour statistics} — per-channel (RGB) mean intensity
        histograms over 200 sampled images indicate that AI-generated images
        tend to have slightly higher mean blue channel values than authentic
        photographs, consistent with synthetic sky or background overrepresentation.
  \item \textbf{File size} — file-size histograms and box-plots show that
        AI-generated images are on average 15--25\,KB larger than authentic
        photographs, suggesting higher JPEG quality or PNG usage by generative
        pipelines.
\end{enumerate}

All images are resized to $224 \times 224$ pixels.  During training, a
randomised augmentation pipeline is applied to reduce overfitting:
\begin{itemize}
  \item Random horizontal flip ($p=0.5$) and vertical flip ($p=0.2$)
  \item Random rotation ($\pm 20^{\circ}$) and colour jitter (brightness~0.3,
        contrast~0.3, saturation~0.2, hue~0.1)
  \item Random grayscale ($p=0.1$), random affine (translation~5\%, scale
        0.9--1.1, shear~5$^{\circ}$)
  \item Centre-crop at 95\% of $224$, then resize back to $224 \times 224$
  \item ImageNet normalisation ($\mu=[0.485,0.456,0.406]$,
        $\sigma=[0.229,0.224,0.225]$)
\end{itemize}
Validation and test images receive only resize and normalisation.

\subsection{Backbone Architectures}

\textbf{ConvNeXt V2 Tiny.}  The backbone (\texttt{convnextv2\_tiny.fcmae\_ft\-\_in22k\_in1k},
28.1\,M parameters) is pretrained on ImageNet-22k with FCMAE self-supervised
pretraining and subsequently fine-tuned on ImageNet-1k.  A custom binary head
is appended:
\begin{equation}
  \text{Linear}(768 \to 256) \to \text{GELU} \to \text{Dropout}(0.3)
        \to \text{Linear}(256 \to 1).
  \label{eq:head}
\end{equation}
Training proceeds in three phases with AdamW and cosine annealing: a warmup
phase (last three ConvNeXt stages + head, 5~epochs, lr~$=10^{-4}$); a
fine-tuning phase (same layers, 15~epochs, lr~$=3\times10^{-5}$); and an
optional full-backbone phase (all weights, 10~epochs, lr~$=10^{-5}$) triggered
automatically when validation F1~$<0.90$.  Gradient clipping
(max\_norm~$=1.0$) is applied throughout.

\textbf{Swin Transformer Tiny.}  The backbone (\texttt{swin\_tiny\_patch4\-\_window7\_224},
27.7\,M parameters) processes $4 \times 4$ non-overlapping patches with
shifted-window attention at multiple resolutions.  An identical binary head
(Eq.~\ref{eq:head}) is attached externally.  Because \texttt{timm}~1.0+ returns
spatial tensors of shape $(N, H, W, C)$ from \texttt{forward\_features()}, a
global average pooling step is applied inside the wrapper model before the
classification head, ensuring the head always receives an $(N, 768)$ pooled
feature vector.  Fine-tuning unfreezes the last three Swin stages plus the
layer-normalisation head and trains for 10~epochs (lr~$=5\times10^{-5}$),
with an optional 10-epoch extension at lr~$=10^{-5}$ if validation F1 remains
below 0.90.

\subsection{Embedding Fusion and Meta-Learner}

Forward hooks capture the 256-dimensional GELU activations from the penultimate
layer of each backbone.  For every image, the two 256-d vectors are concatenated
to form a 512-d representation:
\begin{equation}
  \mathbf{z} = [\,\mathbf{e}_{\text{ConvNeXt}}\;\|\;\mathbf{e}_{\text{Swin}}\,]
        \in \mathbb{R}^{512}.
  \label{eq:fusion}
\end{equation}
The feature matrix is standardised (zero mean, unit variance) using a
\texttt{StandardScaler} fitted on training embeddings only.  Three meta-learners
are trained on $X_{\mathrm{train}} \in \mathbb{R}^{3840 \times 512}$: XGBoost
(max\_depth~$=4$, lr~$=0.03$, 300 estimators), CatBoost (500 iterations,
depth~$=6$, lr~$=0.03$, balanced class weights), and TabPFN.  No test embeddings
are used during this step.

\subsection{Test-Time Augmentation and Threshold Optimisation}

Five deterministic image views (original, horizontal flip, $\pm10^{\circ}$
rotation, colour jitter, 90\% centre-crop) are passed through both backbones,
yielding five 512-d feature vectors per sample.  The best meta-learner's
probability outputs are averaged across the five views to obtain a smoothed
probability estimate $\hat{p}$.  A threshold sweep over $\tau \in [0.20, 0.80]$
(61 steps) is conducted on the validation set to maximise F1, and the selected
threshold is applied to produce binary predictions on both validation and test sets.

%% ── 4. Experiments and Results ────────────────────────────────────────────────
\section{Experiments and Results}
\label{sec:experiments}

\subsection{Experimental Setup}

All experiments are run in Python~3.10 with PyTorch (CUDA-enabled), \texttt{timm},
\texttt{catboost}, \texttt{xgboost}, \texttt{pytorch-grad-cam}, and
\texttt{statsmodels} on a Google Colab T4 GPU.  Mixed-precision training (AMP)
is enabled via \texttt{torch.cuda.amp}.  All random seeds are fixed to 42 for
full reproducibility (\texttt{random}, \texttt{numpy}, \texttt{torch},
\texttt{cudnn.deterministic=True}).

\subsection{Backbone Training Results}

\begin{table}[t]
  \caption{Individual Backbone Validation Performance}
  \label{tab:backbones}
  \centering
  \setlength{\tabcolsep}{4pt}
  \begin{tabular}{lcccc}
    \toprule
    \textbf{Model} & \textbf{Params} & \textbf{Epochs} & \textbf{Val F1} & \textbf{Val AUC} \\
    \midrule
    ConvNeXt V2 Tiny & 28.1M & 5+15+10 & 0.8521 & 0.9140 \\
    Swin Tiny        & 27.7M & 10+10   & 0.8378 & 0.9574 \\
    \bottomrule
  \end{tabular}
\end{table}

Table~\ref{tab:backbones} summarises the best validation F1 and ROC-AUC scores
for each backbone after all training phases.  ConvNeXt V2 Tiny achieves a
higher F1 (0.8521 vs.\ 0.8378), while Swin Tiny achieves a notably higher
ROC-AUC (0.9574 vs.\ 0.9140), indicating that Swin's probability rankings are
better calibrated even though its hard-decision accuracy is slightly lower.
Figures showing training/validation loss, F1, and AUC curves across epochs are
produced from notebook cells~2 and~3.

\subsection{Meta-Learner Benchmark}

\begin{table}[t]
  \caption{Meta-Learner Benchmark on 512-d Fused Embeddings}
  \label{tab:metalearners}
  \centering
  \setlength{\tabcolsep}{4pt}
  \begin{tabular}{lccccc}
    \toprule
    \textbf{Model} & \textbf{F1} & \textbf{AUC} & \textbf{PR-AUC} & \textbf{Prec.} & \textbf{Rec.} \\
    \midrule
    CatBoost  & \textbf{0.8639} & \textbf{0.9274} & \textbf{0.9089} & 0.8264 & 0.9050 \\
    XGBoost   & 0.8516          & 0.9229          & 0.8790          & 0.7687 & 0.9546 \\
    TabPFN    & \multicolumn{5}{c}{\textit{Unavailable (API incompatibility)}} \\
    \bottomrule
  \end{tabular}
\end{table}

Table~\ref{tab:metalearners} reports the per-learner metrics after
per-learner threshold optimisation on the validation set.  CatBoost is the top
performer on F1 (0.8639, threshold~$=0.70$), benefiting from its balanced
class-weight setting and early stopping on F1.  XGBoost achieves higher recall
(0.9546) at the cost of lower precision (0.7687), with an optimal threshold of
0.74.  TabPFN was excluded due to an API incompatibility in the execution
environment.  Figure~\ref{fig:metabars} (produced from notebook cell~4) shows
the F1 comparison bar chart.

\begin{figure}[t]
  \centering
  %% Replace with: \includegraphics[width=0.9\columnwidth]{figs/meta_f1_bar.pdf}
  \framebox{\parbox{0.88\columnwidth}{\centering\vspace{1.1cm}%
    [Meta-learner F1 comparison bar chart: CatBoost~0.864, XGBoost~0.852]%
    \vspace{1.1cm}}}
  \caption{Validation F1 comparison across meta-learners.}
  \label{fig:metabars}
\end{figure}

\subsection{TTA and Threshold Optimisation}

Applying five-view TTA with CatBoost and sweeping the decision threshold over
the validation set yields an optimal threshold of $\tau^{*}=0.68$ and a
best validation F1 of \textbf{0.8771}, an improvement of 1.32 percentage points
over the fixed-threshold CatBoost baseline.  Figure~\ref{fig:threshold}
illustrates the threshold vs.\ F1 curve.

\begin{figure}[t]
  \centering
  %% Replace with: \includegraphics[width=\columnwidth]{figs/threshold_f1.pdf}
  \framebox{\parbox{0.92\columnwidth}{\centering\vspace{1.1cm}%
    [Threshold vs.\ Validation F1 curve; peak at $\tau=0.68$, F1$=0.877$]%
    \vspace{1.1cm}}}
  \caption{Validation F1 as a function of decision threshold (CatBoost + 5-view TTA).}
  \label{fig:threshold}
\end{figure}

\subsection{Full Pipeline Evaluation}

The confusion matrix (Figure~\ref{fig:cm}) and ROC curve (AUC displayed from
cell~7 output) characterise the final validation performance.  The
classification report from cell~7 provides per-class precision and recall
figures; false positives (authentic images predicted as AI-generated) are
primarily images with uniform, texture-heavy backgrounds, while false negatives
(AI-generated images missed by the model) tend to display photorealistic
lighting and minimal generative artefacts.  A McNemar's test (using
\texttt{statsmodels}) comparing the ConvNeXt V2 Tiny standalone classifier
to the full ensemble confirms that the improvement in correct classifications is
statistically significant ($p < 0.05$).

\begin{figure}[t]
  \centering
  %% Replace with: \includegraphics[width=0.75\columnwidth]{figs/confusion_matrix.pdf}
  \framebox{\parbox{0.73\columnwidth}{\centering\vspace{1.2cm}%
    [Confusion matrix heatmap (CatBoost + TTA, val set)]%
    \vspace{1.2cm}}}
  \caption{Confusion matrix on the validation set (CatBoost + 5-view TTA,
           threshold~$=0.68$).}
  \label{fig:cm}
\end{figure}

\subsection{Grad-CAM Error Analysis}

Grad-CAM visualisations (cell~8) highlight the spatial regions driving each
prediction.  For true positives, the model attends to areas with telltale
generative artefacts—over-smoothed edges, unnatural colour gradients, and
repeated texture tiles.  For false positives, attention localises to busy
background patterns that superficially resemble AI generation.  For false
negatives, the saliency map is often diffuse, indicating that the model fails
to identify a single discriminative region in highly photorealistic synthetic
images.  Table~\ref{tab:failures} summarises the identified failure modes.

\begin{table}[t]
  \caption{Systematic Failure Modes Identified via Grad-CAM}
  \label{tab:failures}
  \centering
  \begin{tabular}{p{1.3cm}p{5.2cm}}
    \toprule
    \textbf{Type} & \textbf{Root cause} \\
    \midrule
    FP & Misleading background texture resembling generative artefacts \\
    FP & Handcrafted repetitive patterns confusing local texture detector \\
    FN & Severe lighting artefacts washing out discriminative features \\
    FN & Multi-object scenes where model attends to the wrong region \\
    \bottomrule
  \end{tabular}
\end{table}

\subsection{Final Submission}

The pipeline is applied to all 2\,058 unlabelled test images, producing a
balanced prediction distribution: 1\,037 images predicted as AI-generated (1)
and 1\,021 as authentic (0).  The output is saved as
\texttt{dcu\_2026\_final\_submission.csv} with columns matching the image-id
and label columns of \texttt{train.csv}, as required by
\texttt{getting\_started.ipynb}.

%% ── 5. Discussion ─────────────────────────────────────────────────────────────
\section{Discussion}
\label{sec:discussion}

The results confirm that the dual-backbone fusion paradigm is more effective
than either backbone used in isolation.  ConvNeXt V2's FCMAE pretraining
provides strong local texture representations, while Swin Transformer's
shifted-window attention captures global context; together they offer
complementary information that the meta-learner exploits to lift F1 by
approximately 2.5 percentage points over the best individual backbone.

The significant gap between training loss (which converges to ${\approx}0.01$)
and validation loss (which rises beyond 1.0 in later fine-tuning epochs) suggests
residual overfitting despite aggressive augmentation.  Larger pre-training
datasets, stronger regularisation (e.g., DropPath, stochastic depth), or
self-supervised fine-tuning on unlabelled listing images could address this.

A practical limitation is that the pipeline's total inference time scales with
the number of TTA views and two heavy backbones, making real-time deployment
at Etsy's scale non-trivial.  Model distillation or caching of backbone
embeddings could mitigate latency.

%% ── 6. Conclusion ─────────────────────────────────────────────────────────────
\section{Conclusion}
\label{sec:conclusion}

We presented an end-to-end pipeline for detecting AI-generated product listing
images on Etsy, combining ConvNeXt V2 Tiny and Swin Transformer Tiny with
a CatBoost meta-learner on fused 512-dimensional embeddings.  Five-view
test-time augmentation and validation-set threshold optimisation improve the
validation F1 to \textbf{0.8771}.  Grad-CAM analysis provides interpretable
insight into model failures.  Future work will investigate frequency-domain
features, CLIP-based universal detection~\cite{ojha2023towards}, and knowledge
distillation for production deployment.

%% ── Team Member Contributions ─────────────────────────────────────────────────
\section*{Team Member Contributions}

\textbf{Indumanaswini Bhukya:} Led the backbone training and evaluation components,
including the ConvNeXt V2 Tiny and Swin Transformer Tiny fine-tuning schedules,
gradient clipping strategy, and checkpoint selection.  Implemented the embedding
fusion pipeline, the 512-d feature extraction via forward hooks, and the
StandardScaler normalisation.  Conducted Grad-CAM visualisation and failure-mode
analysis.  Wrote Sections~I, III (Backbones, Embedding Fusion), IV.B, IV.D, and~V.

\textbf{Rajesh Suddala:} Led the meta-learner benchmarking and inference pipeline,
including XGBoost and CatBoost training, per-learner threshold optimisation, and
the five-view test-time augmentation implementation.  Managed the dataset loading,
stratified splitting, and data augmentation pipeline, and verified the integrity of
the final submission.  Wrote Sections~II, III (TTA, Threshold Optimisation),
IV.A, IV.C, and~VI.

%% ── References ────────────────────────────────────────────────────────────────
\bibliographystyle{IEEEtran}
\begin{thebibliography}{99}

\bibitem{ramesh2022dalle}
A.~Ramesh, P.~Dhariwal, A.~Nichol, C.~Chu, and M.~Chen,
``Hierarchical text-conditional image generation with CLIP latents,''
\textit{arXiv preprint arXiv:2204.06125}, 2022.

\bibitem{rombach2022latent}
R.~Rombach, A.~Blattmann, D.~Lorenz, P.~Esser, and B.~Ommer,
``High-resolution image synthesis with latent diffusion models,''
in \textit{Proc. IEEE/CVF CVPR}, pp.~10684--10695, 2022.

\bibitem{woo2023convnextv2}
S.~Woo, S.~Debnath, R.~Hu, X.~Chen, Z.~Liu, I.~S.~Kweon, and S.~Xie,
``ConvNeXt V2: Co-designing and scaling ConvNets with masked autoencoders,''
in \textit{Proc. IEEE/CVF CVPR}, pp.~16133--16142, 2023.

\bibitem{liu2021swin}
Z.~Liu, Y.~Lin, Y.~Cao, H.~Hu, Y.~Wei, Z.~Zhang, S.~Lin, and B.~Guo,
``Swin Transformer: Hierarchical vision transformer using shifted windows,''
in \textit{Proc. IEEE/CVF ICCV}, pp.~10012--10022, 2021.

\bibitem{chen2016xgboost}
T.~Chen and C.~Guestrin,
``XGBoost: A scalable tree boosting system,''
in \textit{Proc. ACM KDD}, pp.~785--794, 2016.

\bibitem{prokhorenkova2018catboost}
L.~Prokhorenkova, G.~Gusev, A.~Vorobev, A.~V.~Dorogush, and A.~Gulin,
``CatBoost: Unbiased boosting with categorical features,''
in \textit{Proc. NeurIPS}, pp.~6638--6648, 2018.

\bibitem{selvaraju2017gradcam}
R.~R.~Selvaraju, M.~Cogswell, A.~Das, R.~Vedantam, D.~Parikh, and D.~Batra,
``Grad-CAM: Visual explanations from deep networks via gradient-based
localization,''
in \textit{Proc. IEEE/CVF ICCV}, pp.~618--626, 2017.

\bibitem{wang2020cnn}
S.-Y.~Wang, O.~Wang, R.~Zhang, A.~Owens, and A.~A.~Efros,
``CNN-generated images are surprisingly easy to spot \ldots for now,''
in \textit{Proc. IEEE/CVF CVPR}, pp.~8695--8704, 2020.

\bibitem{frank2020leveraging}
J.~Frank, T.~Eisenhofer, L.~Schiele, A.~Fischer, D.~Kolossa, and T.~Holz,
``Leveraging frequency analysis for deep fake image recognition,''
in \textit{Proc. ICML}, pp.~3247--3258, 2020.

\bibitem{corvi2023detection}
R.~Corvi, D.~Cozzolino, G.~Zingarini, G.~Poggi, K.~Nagano, and L.~Verdoliva,
``On the detection of synthetic images generated by diffusion models,''
in \textit{Proc. IEEE ICASSP}, pp.~1--5, 2023.

\bibitem{ojha2023towards}
U.~Ojha, Y.~Li, and Y.~J.~Lee,
``Towards universal fake image detection leveraging vision foundation models,''
in \textit{Proc. IEEE/CVF CVPR}, pp.~24205--24214, 2023.

\bibitem{radford2021learning}
A.~Radford \textit{et al.},
``Learning transferable visual models from natural language supervision,''
in \textit{Proc. ICML}, pp.~8748--8763, 2021.

\bibitem{farid2009image}
H.~Farid,
``Image forgery detection,''
\textit{IEEE Signal Process. Mag.}, vol.~26, no.~2, pp.~16--25, Mar. 2009.

\bibitem{provost2000wellbehaved}
F.~Provost and T.~Fawcett,
``Well-trained PETs: Improving probability estimation trees,''
AAAI Technical Report, 2000.

\end{thebibliography}

\end{document}
